\documentclass[10pt,a4paper]{amsart}
\usepackage{amsmath,amssymb,amsthm,graphicx}
% Retained preamble for the Stacks Project chapter snapshot in this directory.
% The chapter's own preamble.tex is a copy of the shared Stacks Project preamble; the block
% below is sliced MECHANICALLY, line for line, from it. Nothing was added, removed or
% redefined; the wrapper class, the xr-hyper/hyperref wiring and the externaldocument list
% are omitted because the assembler supplies the document class.
%
% lines taken from preamble.tex: 3, 5, 9, 11, 12, 156, 157, 158, ... (51 lines total)
%
% package coverage restored from the chapter's own preamble (2026-10-13): xy, xy-2cell

\usepackage[all]{xy}\usepackage{verbatim}
\xyoption{2cell}
\UseAllTwocells

\newenvironment{reference}{\comment}{\endcomment}
\newenvironment{slogan}{\comment}{\endcomment}
\newenvironment{history}{\comment}{\endcomment}
\theoremstyle{plain}
\newtheorem{theorem}[subsection]{Theorem}
\newtheorem{proposition}[subsection]{Proposition}
\newtheorem{lemma}[subsection]{Lemma}
\theoremstyle{definition}
\newtheorem{definition}[subsection]{Definition}
\newtheorem{example}[subsection]{Example}
\newtheorem{exercise}[subsection]{Exercise}
\newtheorem{situation}[subsection]{Situation}
\theoremstyle{remark}
\newtheorem{remark}[subsection]{Remark}
\newtheorem{remarks}[subsection]{Remarks}
\numberwithin{equation}{subsection}
\def\lim{\mathop{\mathrm{lim}}\nolimits}
\def\colim{\mathop{\mathrm{colim}}\nolimits}
\def\Spec{\mathop{\mathrm{Spec}}}
\def\Hom{\mathop{\mathrm{Hom}}\nolimits}
\def\Ext{\mathop{\mathrm{Ext}}\nolimits}
\def\SheafHom{\mathop{\mathcal{H}\!\mathit{om}}\nolimits}
\def\SheafExt{\mathop{\mathcal{E}\!\mathit{xt}}\nolimits}
\def\Sch{\mathit{Sch}}
\def\Mor{\mathop{\mathrm{Mor}}\nolimits}
\def\Ob{\mathop{\mathrm{Ob}}\nolimits}
\def\Sh{\mathop{\mathit{Sh}}\nolimits}
\def\NL{\mathop{N\!L}\nolimits}
\def\CH{\mathop{\mathrm{CH}}\nolimits}
\def\proetale{{pro\text{-}\acute{e}tale}}
\def\etale{{\acute{e}tale}}
\def\QCoh{\mathit{QCoh}}
\def\Ker{\mathop{\mathrm{Ker}}}
\def\Im{\mathop{\mathrm{Im}}}
\def\Coker{\mathop{\mathrm{Coker}}}
\def\Coim{\mathop{\mathrm{Coim}}}
\DeclareMathSymbol{\boxtimes}{\mathbin}{AMSa}{"02}
\def\QCohstack{\mathcal{QC}\!\mathit{oh}}
\def\Cohstack{\mathcal{C}\!\mathit{oh}}
\def\Spacesstack{\mathcal{S}\!\mathit{paces}}
\def\Quotfunctor{\mathrm{Quot}}
\def\Hilbfunctor{\mathrm{Hilb}}
\def\Curvesstack{\mathcal{C}\!\mathit{urves}}
\def\Polarizedstack{\mathcal{P}\!\mathit{olarized}}
\def\Complexesstack{\mathcal{C}\!\mathit{omplexes}}
\def\Pic{\mathop{\mathrm{Pic}}\nolimits}
\def\Picardstack{\mathcal{P}\!\mathit{ic}}
\def\Picardfunctor{\mathrm{Pic}}
\def\Deformationcategory{\mathcal{D}\!\mathit{ef}}
\usepackage{cleveref}
\crefname{originallemma}{original lemma}{original lemmas}
\Crefname{originallemma}{Original lemma}{Original lemmas}
\makeatletter
\newlabel{lemma-quasi-coherence-direct-image}{{08FA}{source}{}{sourcepointer-1}{}}
\newlabel{lemma-quasi-coherence-direct-image@cref}{{[originallemma][0][]08FA}{[1][0][]source}{}{}{}}
\makeatother
\crefname{originallemma}{original lemma}{original lemmas}
\Crefname{originallemma}{Original lemma}{Original lemmas}
\makeatletter
\newlabel{theorem-approximation}{{08HP}{source}{}{sourcepointer-2}{}}
\newlabel{theorem-approximation@cref}{{[originallemma][0][]08HP}{[1][0][]source}{}{}{}}
\makeatother
\title{Stacks, Derived Categories of Spaces: perfect-enough}
\author{The Stacks Project Authors}
\date{stacks-project:spaces-perfect.tex:local-snapshot-20261012; GNU FDL 1.2 or later}
\begin{document}
\maketitle
\noindent\textit{Editorial note.} Statement span is lines 6543-6551 of sources/stacks-spaces-perfect/pinned-original-source.tex; proof span is lines 6553-6581, the complete author proof (27 body lines).\par\medskip
\noindent\textit{Editorial note.} Navigation stubs: 2 cited-but-not-included labels are supplied as non-load-bearing pointers with their pinned tags where the snapshot tag index has them.\par\medskip
\noindent\textit{Editorial note.} extra\_packages is empty on purpose: the retained preamble already loads xy with options and a second declaration clashes in the assembler wrapper.\par\medskip
\noindent\textit{Editorial note.} Text is retained verbatim, including the author's own typos and placeholders.\par\medskip
\noindent\textit{Editorial note.} Queue index 41 of queue\_spaceperfect.json (pb=1146, stmt 6543-6551, proof 6553-6581).\par\medskip
\noindent\textit{Editorial note.} stacks-spaces-perfect wave sp01 (2026-10-14): fresh adjudication accept per lane-s3/sp01\_judgments; pool per Chief 12:4x unfreeze (reserve-pool workflow, six-check clean); spaces-perfect.tex 126 with-proof rows, 70 b107-occupied labels avoided, 56 free rows judged 32 accept / 24 reject.\par\medskip
\noindent\textit{Editorial note.} Cross-chapter same-label warning (NOT an occupancy conflict, Chief ruling 2026-10-13 evening): label 'lemma-perfect-enough' also exists in stacks-perfect — a different proposition in a different chapter file.\par\medskip
\noindent\textit{Allowed standard prerequisites (editorial).}\begin{itemize}
\item Explicit prerequisite boundary. The cited inputs are prerequisites; the argument in the retained span is not replaced by any cited result.
\end{itemize}\medskip
\ifdefined\hypertarget\hypertarget{sourcepointer-1}{}\fi
\noindent\textit{Source reference.} Original lemma 08FA: sources/stacks-spaces-perfect/pinned-original-source.tex, source lines 790--804. Cited by the retained statement or proof but not part of the collected spans; supplied only as a navigation pointer. The argument does not depend on its content.\par\medskip
\ifdefined\hypertarget\hypertarget{sourcepointer-2}{}\fi
\noindent\textit{Source reference.} Original lemma 08HP: sources/stacks-spaces-perfect/pinned-original-source.tex, source lines 3786--3790. Cited by the retained statement or proof but not part of the collected spans; supplied only as a navigation pointer. The argument does not depend on its content.\par\medskip
\begin{lemma}
\label{lemma-perfect-enough}
Let $A$ be a ring. Let $X$ be an algebraic space over $A$ which is
quasi-compact and quasi-separated. Let $K \in D^-_\QCoh(\mathcal{O}_X)$.
If $R\Gamma(X, E \otimes^\mathbf{L} K)$ is pseudo-coherent
in $D(A)$ for every perfect $E$ in $D(\mathcal{O}_X)$,
then $R\Gamma(X, E \otimes^\mathbf{L} K)$ is pseudo-coherent
in $D(A)$ for every pseudo-coherent $E$ in $D(\mathcal{O}_X)$.
\end{lemma}

\begin{proof}
There exists an integer $N$ such that
$R\Gamma(X, -) : D_\QCoh(\mathcal{O}_X) \to D(A)$
has cohomological dimension $N$ as explained in
Lemma \ref{lemma-quasi-coherence-direct-image}.
Let $b \in \mathbf{Z}$ be such that $H^i(K) = 0$ for $i > b$.
Let $E$ be pseudo-coherent on $X$.
It suffices to show that $R\Gamma(X, E \otimes^\mathbf{L} K)$
is $m$-pseudo-coherent for every $m$.
Choose an approximation $P \to E$ by a perfect complex $P$
of $(X, E, m - N - 1 - b)$. This is possible by
Theorem \ref{theorem-approximation}.
Choose a distinguished triangle
$$
P \to E \to C \to P[1]
$$
in $D_\QCoh(\mathcal{O}_X)$. The cohomology sheaves of $C$ are zero
in degrees $\geq m - N - 1 - b$. Hence the cohomology
sheaves of $C \otimes^\mathbf{L} K$ are zero in degrees $\geq m - N - 1$.
Thus the cohomology of $R\Gamma(X, C \otimes^\mathbf{L} K)$
are zero in degrees $\geq m - 1$. Hence
$$
R\Gamma(X, P \otimes^\mathbf{L} K) \to R\Gamma(X, E \otimes^\mathbf{L} K)
$$
is an isomorphism on cohomology in degrees $\geq m$.
By assumption the source is pseudo-coherent.
We conclude that $R\Gamma(X, E \otimes^\mathbf{L} K)$
is $m$-pseudo-coherent as desired.
\end{proof}

\end{document}
